% ID: SEL-D0071-OA-E03
% ESTADO: cotejada. Segunda pasada visual y borrador compilado revisados.
% Fuente: Selectividad/sel_2009_matematicas_aplicadas/mat apli CC SS05.pdf
% SHA-256: ca51ab435c7584a8e047e19d0826f3273f3e88aed6110ccd00e46a40e85518b1
% Página: 1; opción A; ejercicio original 3.

\noindent\underline{Parte I}

Sean $A$ y $B$ dos sucesos tales que $P(A)=0.3$, $P(B)=0.4$, $P(A\cup B)=0.65$. Conteste razonadamente las siguientes preguntas:
\begin{enumerate}[label=\alph*)]
\item \textbf{(0.5 puntos)} \textquestiondown{}Son incompatibles $A$ y $B$?
\item \textbf{(0.5 puntos)} \textquestiondown{}Son independientes $A$ y $B$?
\item \textbf{(1 punto)} Calcule $P(A/B^C)$.
\end{enumerate}
\par\noindent\begin{minipage}{\linewidth}
\noindent\underline{Parte II}

Una variable aleatoria $X$ se distribuye de forma Normal, con media $\mu$ y desviación típica $\sigma=0.9$.
\begin{enumerate}[label=\alph*)]
\item \textbf{(1 punto)} Una muestra aleatoria de tamaño 9 ha proporcionado los siguientes valores de $X$:
\[
7.0,\ 6.4,\ 8.0,\ 7.1,\ 7.3,\ 7.4,\ 5.6,\ 8.8,\ 7.2.
\]
Obtenga un intervalo de confianza para la media $\mu$, con un nivel de confianza del $97\%$.
\item \textbf{(1 punto)} Con otra muestra, se ha obtenido que un intervalo de confianza para $\mu$, al $95\%$, es el siguiente $(6.906,\ 7.494)$. \textquestiondown{}Cuál es el tamaño de la muestra utilizada?
\end{enumerate}
\end{minipage}\par
